\documentclass[11pt,a4paper]{article}

\usepackage[a4paper,margin=1.7cm]{geometry}
\usepackage[T1]{fontenc}
\usepackage{lmodern}
\usepackage{microtype}
\usepackage{xcolor}
\usepackage{tabularx}
\usepackage{array}
\usepackage{enumitem}
\usepackage{amssymb}
\usepackage[hidelinks]{hyperref}
\usepackage[most]{tcolorbox}

\definecolor{sgsNavy}{HTML}{0E214B}
\definecolor{sgsGold}{HTML}{C7A34A}
\definecolor{sgsGrey}{HTML}{6B7280}
\definecolor{sgsSoft}{HTML}{F4F6FA}

\setlength{\parindent}{0pt}
\setlength{\parskip}{0.5em}

\newtcolorbox{SGSHeaderBox}{
  colback=sgsNavy,
  colframe=sgsNavy,
  boxrule=0pt,
  arc=8pt,
  left=10pt,
  right=10pt,
  top=10pt,
  bottom=10pt
}

\newtcolorbox{SGSCallout}[1]{
  colback=sgsSoft,
  colframe=sgsGold,
  boxrule=0.6pt,
  arc=8pt,
  left=10pt,
  right=10pt,
  top=10pt,
  bottom=10pt,
  title=\textbf{#1},
  fonttitle=\color{sgsNavy}
}

\newtcolorbox{SGSStep}[2]{
  colback=white,
  colframe=sgsGold,
  boxrule=0.6pt,
  arc=10pt,
  left=10pt,
  right=10pt,
  top=10pt,
  bottom=10pt,
  title=\textbf{Step #1: #2},
  fonttitle=\color{sgsNavy}
}

\newcommand{\NameLine}{\textbf{Name:} \rule{7cm}{0.4pt}\hfill \textbf{Class:} \rule{3cm}{0.4pt}\hfill \textbf{Date:} \rule{3cm}{0.4pt}}
\begin{document}

\begin{SGSHeaderBox}
  {\color{white}\LARGE\textbf{Unit 2 Activity 1: Who is John Doe?}}\\
  {\color{white}\large Investigating digital footprints}
\end{SGSHeaderBox}

\vspace{0.6em}
\NameLine

\begin{SGSCallout}{What you are learning}
You can build a detailed picture of someone from small pieces of public information. This activity helps you spot what a digital footprint reveals and decide what should stay private.
\end{SGSCallout}

\begin{SGSStep}{1}{Open the John Doe Dossier}
\textbf{Open the dossier on the website:} \url{https://sgs-science.web.app/ks3/year7/unit2/john-doe.html}

\begin{itemize}[leftmargin=*,label=$\square$]
  \item Scan all 15 evidence images.
  \item For each image, ask: \textit{What does this reveal about John?}
\end{itemize}
\end{SGSStep}

\begin{SGSStep}{2}{Complete the Investigation Table}
Fill in one row for each evidence image.

\vspace{0.4em}
\renewcommand{\arraystretch}{1.2}
\begin{tabularx}{\textwidth}{|p{1.7cm}|X|p{4.4cm}|}
\hline
\textbf{Image \#} & \textbf{What does it reveal?} & \textbf{Public or private? Why?} \\
\hline
1 & \rule{0pt}{1.2cm} & \rule{0pt}{1.2cm} \\ \hline
2 & \rule{0pt}{1.2cm} & \rule{0pt}{1.2cm} \\ \hline
3 & \rule{0pt}{1.2cm} & \rule{0pt}{1.2cm} \\ \hline
4 & \rule{0pt}{1.2cm} & \rule{0pt}{1.2cm} \\ \hline
5 & \rule{0pt}{1.2cm} & \rule{0pt}{1.2cm} \\ \hline
6 & \rule{0pt}{1.2cm} & \rule{0pt}{1.2cm} \\ \hline
7 & \rule{0pt}{1.2cm} & \rule{0pt}{1.2cm} \\ \hline
8 & \rule{0pt}{1.2cm} & \rule{0pt}{1.2cm} \\ \hline
9 & \rule{0pt}{1.2cm} & \rule{0pt}{1.2cm} \\ \hline
10 & \rule{0pt}{1.2cm} & \rule{0pt}{1.2cm} \\ \hline
11 & \rule{0pt}{1.2cm} & \rule{0pt}{1.2cm} \\ \hline
12 & \rule{0pt}{1.2cm} & \rule{0pt}{1.2cm} \\ \hline
13 & \rule{0pt}{1.2cm} & \rule{0pt}{1.2cm} \\ \hline
14 & \rule{0pt}{1.2cm} & \rule{0pt}{1.2cm} \\ \hline
15 & \rule{0pt}{1.2cm} & \rule{0pt}{1.2cm} \\ \hline
\end{tabularx}
\end{SGSStep}

\begin{SGSStep}{3}{Find 10+ Clues}
Using the evidence, find at least \textbf{10 pieces of information} about John. Your list must include:
\begin{itemize}[leftmargin=*,label=$\square$]
  \item John's name
  \item John's date of birth (day + month)
  \item The city John lives in
\end{itemize}
After that, add any other clues you can work out (hobbies, routine, places he visits, things he likes, family information, etc.).
\end{SGSStep}

\begin{SGSStep}{4}{Decide: Public or Private?}
For each clue in your table, explain whether it should be public or private.
\begin{itemize}[leftmargin=*,label=$\square$]
  \item Could this help someone identify or locate John?
  \item Could this harm John's reputation or safety?
  \item Is it \textit{nice to share} or \textit{risky to share}?
\end{itemize}
\end{SGSStep}

\begin{SGSStep}{5}{Reflection}
\textbf{A. What picture can you build of John Doe?}\\
Write 4--5 bullet points describing what kind of person John appears to be (based only on the evidence).

\vspace{0.4em}
\rule{\textwidth}{0.4pt}\\
\rule{\textwidth}{0.4pt}\\
\rule{\textwidth}{0.4pt}\\
\rule{\textwidth}{0.4pt}\\
\rule{\textwidth}{0.4pt}\\

\vspace{0.8em}
\textbf{B. What should have stayed private?}\\
Choose three items from your table that you think should not have been shared publicly. For each one, write one sentence explaining the risk or concern.

\vspace{0.4em}
\rule{\textwidth}{0.4pt}\\
\rule{\textwidth}{0.4pt}\\
\rule{\textwidth}{0.4pt}\\
\end{SGSStep}

\begin{SGSStep}{6}{Partner Discussion}
Discuss your findings with the person next to you. Share one thing you discovered that you didn't expect.
\end{SGSStep}

\begin{SGSCallout}{By the end of this activity, you should be able to...}
\begin{itemize}[leftmargin=*,label=$\square$]
  \item Explain the difference between public and private information online
  \item Recognise how small pieces of data can build a picture of someone
  \item Understand why sharing too much online can be risky
  \item Describe what responsible digital behaviour looks like
\end{itemize}
\end{SGSCallout}

\vfill
{\small\color{sgsGrey}Printable version generated for SGS Science. Online version: \url{https://sgs-science.web.app/ks3/year7/unit2/activity1.html}}

\end{document}
